\documentclass[a4paper,12pt]{article}

\usepackage[T2A]{fontenc}
\usepackage[utf8]{inputenc}
\usepackage[russian]{babel}

\usepackage{amsmath, amssymb}
\usepackage{geometry}
\usepackage{fancyhdr}
\usepackage{enumitem}

\geometry{margin=2cm}
\setlength{\parskip}{6pt}
\setlength{\parindent}{0pt}

\pagestyle{fancy}
\fancyhf{}
\cfoot{\thepage}
\rhead{8 класс}

\begin{document}

\begin{center}
    {\Large \textbf{Делимость. Теорема Вильсона.}}\\[0.5em]
\end{center}
\noindent Мотивация: хотим что-то знать про сравнение по простому модулю. На самом деле вычеты по простому модулю образуют поле.

\noindent План доказательства: \\
1. Узнаем, что такое поле, кольцо. \\
2. Осознаем понятие класса эквивалентности, вычета. \\
3. Узнаем, что такое ПСВ, свойство ПСВ. \\
4. Аналогично для ПрСВ. \\
5. Узнаем, что такое кольцо вычетов и какие элементы в нем обратимы. \\
6. Кольцо вычетов по простому модулю - поле. \\
7. Следствие: много чего, в частности - теорема Вильсона.

    
\begin{enumerate}[label=\textbf{Задача \arabic*.}]
\item Найдите все решения сравнений: \\ 
а) $4x \equiv 9 \pmod{13}$; \\ 
б) $3x \equiv 12 \pmod{15}$; \\
в) $20x \equiv 30 \pmod{55}$. \\
\item Решите систему сравнений $\begin{cases} 6x + 5y \equiv 1 \pmod{11}, \\ 4x + 3y \equiv 2 \pmod{11}. \end{cases}$
\item Докажите, что $(p - 2)! \equiv 1 \pmod{p}$ при любом простом $p$. \\
\item Докажите, что если $p = 4k + 1$ — простое число, то $x = \left(\frac{p-1}{2} \right)!$ удовлетворяет сравнению $x^2 \equiv -1 \pmod{p}$. \\
\item С помощью теоремы Вильсона докажите, что среди чисел вида $n! + 1$ бесконечно много составных.
\item Пусть $p(x)$ — многочлен с целыми коэффициентами и $a \equiv b \pmod{m}$. Тогда $p(a) \equiv p(b) \pmod{m}$.
\item Докажите, что среди любых n последовательных натуральных чисел есть число, которое делится нацело на n.
\item Докажите, что число является квадратом целого числа тогда и только тогда, когда оно имеет нечетное число натуральных делителей.
\item Верно ли, что бесконечных чисел бесконечно много?
\item Докажите, что простых чисел вида $6k - 1$ бесконечно много.
\item Натуральное число n не имеет собственных делителей, больших 1 и не превосходящих $\sqrt{n}$. Докажите, что n - простое число.
\end{enumerate}
\end{document}